\maketitle%与上面的\title对应
\thispagestyle{plain}%摘要专用页显示页码 1，并与后续正文连续编号
\begin{abstract}%摘要+摘要+摘要+摘要+摘要

	山区洪涝灾害常导致道路中断、通信基站损毁，无人机成为打通“最后一公里”应急物资投送与通信保障的关键装备。本文围绕\textbf{临时调度中心 O01 与 15 个服务区、80 个不可拆货箱、8 架异构运输无人机、2 架中继无人机}构成的协同任务，建立了一套\textbf{物理口径统一、四问层层递进、结果可独立复算}的模型体系，并对若干违反直觉的结论给出了严格论证。

	\textbf{对于问题一}，首先由 30\,m 数字高程模型（DEM）沿航段采样最高地面高程，按“最高点 $+50$\,m”确定巡航海拔，逐段求得 240 条航段的水平距离、爬升与下降高度。据此建立\textbf{载荷—航程关系} $L_{g}(q)=L_{g0}-(L_{g0}-L_{gF})(q/Q_{g})^{3/2}$ 与\textbf{逐段能耗模型}：水平巡航能耗由航程反推 $E^{hor}=(d/L_{g}(q))E_{g}^{use}$，爬升附加能耗按效率折算 $E^{up}=mgh^{+}/\eta_{up}$，下降能耗效率为 0 故不单独计。最大安全载荷由能量约束取等号定义，因指数 $3/2$ 无初等反函数，用\textbf{单调性 $+$ Brent 法数值反解}，回代相对误差 $5.6\times10^{-15}$。结果表明\textbf{结构载重而非能量才是主导约束}：A 型 15/15 个服务区全部受 $Q_{g}$ 约束，B 型 14/15，仅 C 型在 5 个远距离服务区受能量约束（如 S008 反解得 58.99\,kg）。随后把每个服务区化为\textbf{质量—体积二维装箱}问题用 FFD 求解，并推导 Martello--Toth 型解析下界。最终得到\textbf{18 个架次（全部 C 型）}的组批方案，\textbf{逐服务区等于解析下界，架次数维度可证最优}；总运输能耗 \textbf{75.07\,kWh}，累计作业 9.37\,h。$\rho_{g}$ 敏感性显示：$\rho_{g}$ 由 0.20 增至 0.35 使架次由 \textbf{18 增至 25}、能耗升至 106.60\,kWh；$\rho_{g}\ge0.375$ 时部分远距离服务区\textbf{即使空载也无法返航，问题开始无可行解}。

	\textbf{对于问题二}，将单点往返扩展为\textbf{异构多点多架次调度}，联合决策货箱组批、访问顺序、机型、实体无人机、共享电池与各架次开始时刻。多点串飞的载荷\textbf{沿航段递减}、能耗必须逐段累加，交接时间按 $\sum_{\text{站}}(t_{h0}+k_s t_{h1})$ 逐站累加。采用“构造—选序—局部搜索—时段驱动贪心派发”四阶段求解，并以两阶段充电模型刻画电池周转；其中首批箱按机队槽位\textbf{轮转播种}（4 架 A $\to$ 2 架 B $\to$ 2 架 C），派发平局按\textbf{最小松弛量} $slack=\min_{b}(d_{b}-t_{b}^{deliver})$ 而非能耗决出。得到\textbf{25 个架次（A 型 8 $+$ B 型 10 $+$ C 型 7，8 架无人机全部投入使用）}、总能耗 \textbf{77.30\,kWh}、完工 \textbf{3.02\,h}（10855.1\,s）、期望送达准时率 \textbf{100.0\%（80/80 箱全部按时送达）}，首批违规与期望违规均为 \textbf{0 箱}。压缩架次数的正确做法不是调高架次权重：单纯的“权重法”虽能把架次数压到 25，却会让真实排程冒出 4 箱超期（准时率降至 95.0\%）而目标函数看不见——把每违规箱权重取 $1$ 与取 $10^{3}$ 得到逐位相同的结果，证明该目标对真实违规是盲的；本文改为以“\textbf{真实调度器判定 0 违规}”为硬前提、以“架次数最少”为唯一方向，对每一次两两合并（必要时升级机型）与整装重排\textbf{都过一遍真实调度器}、不通过即回退，并修掉局部搜索 relocate 走法缺失 frozen 保护、会把货箱搬出首批专架次的缺陷。为验证方案，本文用 240 条航段缓存对全部架次做了\textbf{逐段独立复算}：能耗与方案上报值一致（总能耗均为 77.3016\,kWh）、全部架次满足能量预算，资源占用区间\textbf{零冲突}、结束时刻自洽偏差最大 0.078\,s；独立校验器报出的违规为 \textbf{0 条}（含时限类）。四目标权衡对比表明：同构 C 型方案架次数最少（14），但完工 5.32\,h、准时率仅 62.5\%；本文“机队摊平”方案架次数多 11 个、能耗反而略低（77.30 vs 77.94\,kWh），换来完工时间下降 43\%、准时率 \textbf{100.0\%}；25 个架次距理论下界 24（单轮机队容量 320\,kg / 0.886\,m$^3$ 对比总需求 758\,kg / 2.011\,m$^3$，至少 3 轮）仅高 1 个，质量装载率由 71\% 提升至 \textbf{79\%}（$758/960$）。需要说明的是，本文早期版本曾因\textbf{机队坍缩}（只评估同构机队、把 8 架异构机队压成 2 架 C 型在用）与\textbf{派发退化}（无违规时按能耗打破平局，使首波 8 个名额中 2 个被第二趟 S001 占用）而误判\textbf{“首批保障时限物理不可行”}；修正后可见 60\,min 档服务区恰为 8 个、与 8 架无人机一一对应，\textbf{供给等于需求、恰好可行}——60\,min 档最晚交付 3546.3\,s（S001，裕度 53.7\,s）、120\,min 档最晚 7164.5\,s（S003，裕度 35.5\,s）、180\,min 档最晚 7966.6\,s（S009）。

	\textbf{对于问题三}，新增\textbf{连续通信约束}：运输无人机在爬升、巡航、下降与投送\textbf{全程}不得中断。先对问题二方案做直连诊断，沿完整轨迹以 $\Delta t=2$\,s 逐时刻采样并做\textbf{直连 $>$ 中继 $>$ 中断}三态判定，发现 25 个运输架次中 \textbf{20 个存在直连中断}（占 80.0\%、\textbf{平均中断占比 28.4\%}），另 5 个架次全程直连可用，\textbf{中继无人机因此是必需项}。链路预算给出三条链路的双向门限：直连 122\,dB、中继接入 116\,dB、中继回传 126\,dB，接入段最弱，故中继应\textbf{贴近作业空域而非贴近网关}。将“连续轨迹 $\times$ 连续选址”降维为可验证的充分条件——存在单点 $P$ 使轨迹每一时刻满足“直连 $\vee$（接入 $\wedge$ 回传）”，在 400\,m 网格的 1885 个候选点上择优。最终给出\textbf{20 个中继架次}，需保障的 \textbf{20 个架次全部由单一悬停点实现 100\%（20/20）全程连续通信覆盖}；运输能耗 77.30\,kWh $+$ 中继能耗 4.83\,kWh $=$ \textbf{82.14\,kWh}（中继仅占 5.9\%），联合完工 \textbf{3.15\,h}（11323.3\,s），独立校验器报出 \textbf{0 条违规}。

	\textbf{对于问题四}，把“同一架次涉及的服务区必须划入同一任务组”这一硬规则形式化为服务区\textbf{图上的等价关系}，任务分区即求\textbf{连通分量（原子单元）}。对问题三的 \textbf{25 个架次}构图，其中 \textbf{12 个多点架次把 15 个服务区串成 3 个连通分量（原子单元）}——其中两个单元各只含 1 个服务区（S010、S014），第三个单元含其余 \textbf{13 个服务区}。因此\textbf{合法组数 $K$ 受原子单元拓扑支配}：$K=2$ 与 $K=3$ \textbf{均可由原子单元直接合并得到、改动 0 个架次}。本文进一步定位出 \textbf{6 个桥接架次}（T019、T017、T018、T021、T024、T025）。资源核算表明 K$\,=\,$1/2/3 的四类资源总量为 \textbf{32/34/37}（台$\cdot$组）、组间不均衡度为 \textbf{0/1.8517/2.7035}、最大组工作量由 13.95\,h 降至 13.00\,h，即\textbf{分区不节省资源（总量随 $K$ 单调上升）却能降低单组峰值工作量}；即使 $K=1$ 也已出现缺口（A、B、C 三型共享电池各缺 1 组），随分组增加缺口进一步扩大。

	本文的主要特色在于：\textbf{①物理口径唯一}——附录给定的时间、能耗、充电、链路公式只在公共层实现一次，四问共用，杜绝了口径分叉；\textbf{②结果可独立复算}——用 240 条航段缓存对问题二全部 25 个架次逐段重算，与求解器上报值一致，独立校验器确认问题二、问题三的\textbf{违规均为 0 条}；\textbf{③敢于纠正否定性结论}——早期版本因“机队坍缩 $+$ 派发退化”两处建模缺陷误判“首批保障时限物理不可行”，本文用“8 架无人机 $\leftrightarrow$ 8 个 60\,min 服务区”的\textbf{供给—需求对偶}将其证伪并给出轮转播种与最小松弛量派发；问题四则由精确的图论计算给出“$K=2$、$K=3$ 均可零架次改动分区”的结论。

	\keywords{无人机应急物流；载荷—航程关系；二维装箱；多点多架次调度；连续通信约束；中继选址；图连通分量；可行性校验}
\end{abstract}
